\chapter{汇总头与依赖关系}\label{ch:all}

\chaptermeta{all.hpp}{全部模块}{——}

一次性引入全部模块，并给出各头文件之间的依赖关系与复制顺序。

\section{接口}

\apitable{
    \api{all.hpp}{按文件名顺序包含 \code{def}、\code{modint}、\code{number_theory}、\code{suffix_array}、\code{suffix_automaton}、\code{bipartite_matching}、\code{bipartite_weighted_matching}、\code{general_matching}、\code{general_weighted_matching}、\code{max_flow}、\code{min_cost_flow}、\code{min_cost_circulation}、\code{fpoly}、\code{interpolation}、\code{geometry}、\code{geometry/circle}、\code{geometry/polygon}}
}

\section{约定}

\texttt{all.hpp} 只是把各模块按文件名顺序包含一遍，没有全局初始化逻辑，也不改变任何模块的行为；逐个包含与一次性包含得到的符号完全一致，区别只在编译时间。依赖表中被依赖的文件在前，复制时保持相对路径，不要把 \code{geometry/} 子目录拍平。一次包含全部头文件会实例化更多模板，赛场提交时建议只包含用到的模块。

\subsection{依赖关系}

\apitable{
    \api{def.hpp}{无依赖，所有模块的基础}
    \api{modint.hpp}{def}
    \api{number_theory.hpp}{def}
    \api{suffix_array.hpp}{def}
    \api{suffix_automaton.hpp}{def}
    \api{bipartite_matching.hpp}{def}
    \api{bipartite_weighted_matching.hpp}{def}
    \api{general_matching.hpp}{def}
    \api{general_weighted_matching.hpp}{def}
    \api{max_flow.hpp}{def}
    \api{min_cost_flow.hpp}{def}
    \api{min_cost_circulation.hpp}{def}
    \api{fpoly.hpp}{modint、def}
    \api{interpolation.hpp}{def}
    \api{geometry.hpp}{def}
    \api{geometry/circle.hpp}{geometry、def}
    \api{geometry/polygon.hpp}{geometry、def}
    \api{all.hpp}{以上全部}
}

\subsection{抄写清单}

\begin{itemize}[leftmargin=1.6em, itemsep=0.4em]
    \item 先复制 \code{def.hpp}，它是所有模块的公共依赖。
    \item 按需复制其余模块；只用到几何时，复制 \code{geometry.hpp} 以及 \code{geometry/} 下的子模块。
    \item 头文件之间使用相对路径包含（\code{"def.hpp"}、\code{"../geometry.hpp"}），保持目录结构即可直接编译。
    \item 提交前用 \code{-std=c++20 -O2 -Wall} 编译一次；仓库的 \code{make headers} 会逐个检查头文件能否独立编译。
    \item 若把多个模块合并成单文件提交，注意 \code{#pragma once} 与 \code{namespace acm} 只能出现一次。
\end{itemize}

\fullsource{all.hpp}
